% Part: many-valued-logic
% Chapter: syntax-and-semantics
% Section: sublogics

\documentclass[../../../include/open-logic-section]{subfiles}

\begin{document}

\olfileid[fr]{mvl}{syn}{sub}

\olsection{Les logiques plurivalentes comme
  sous-logiques de~$\LogCL$}

Les logiques plurivalentes usuelles
sont définies par des matrices dont
les fonctions de vérité coïncident
avec les fonctions classiques lorsque
leurs arguments appartiennent à
$\{\True,\False\}$. En particulier,
si $x\in\{\True,\False\}$, alors
$\tf{\lnot}[\Log L](x)
=\tf{\lnot}[\LogCL](x)$.
De même, pour chacun des connecteurs
$\star\in\{\land,\lor,\lif\}$
et tous $x,y\in\{\True,\False\}$,
on a
$\tf{\star}[\Log L](x,y)
=\tf{\star}[\LogCL](x,y)$.
Les restrictions des quatre fonctions
de vérité à ces valeurs sont donc
exactement les fonctions classiques.

\begin{prop}\ollabel{prop:mvl-cl}
Supposons que le langage d'une logique
plurivalente~$\Log L$ contienne
$\lnot$, $\land$, $\lor$ et $\lif$,
que $\True,\False\in V$, et que ses
fonctions de vérité vérifient :
\begin{enumerate}
\item\ollabel{prop:not}
  $\tf{\lnot}[\Log L](x)
   =\tf{\lnot}[\LogCL](x)$
  pour $x\in\{\True,\False\}$;
\item\ollabel{prop:land}
  $\tf{\land}[\Log L](x,y)
   =\tf{\land}[\LogCL](x,y)$;
\item\ollabel{prop:lor}
  $\tf{\lor}[\Log L](x,y)
   =\tf{\lor}[\LogCL](x,y)$;
\item\ollabel{prop:lif}
  $\tf{\lif}[\Log L](x,y)
   =\tf{\lif}[\LogCL](x,y)$,
\end{enumerate}
les trois dernières égalités valant
pour tous $x,y\in\{\True,\False\}$.
Alors, pour toute formule~$!A$
construite uniquement avec des
variables propositionnelles et ces
quatre connecteurs, et pour toute
valuation classique
$\pAssign{v}\colon\PVar\to
\{\True,\False\}$, on a
\[
  \pValue{v}(!A)[\Log L]
    =\pValue{v}(!A)[\LogCL].
\]
\end{prop}

\begin{proof}
Procédons par induction sur~$!A$.
\begin{enumerate}
\item Si $!A\ident p$ est atomique,
  alors
  $\pValue{v}(!A)[\Log L]
   =\pAssign{v}(p)
   =\pValue{v}(!A)[\LogCL]$.
\item Si $!A\ident\lnot !B$, alors
  \begin{align*}
    \pValue{v}(!A)[\Log L]
      &=\tf{\lnot}[\Log L]
        (\pValue{v}(!B)[\Log L])\\
      &=\tf{\lnot}[\Log L]
        (\pValue{v}(!B)[\LogCL])\\
      &=\tf{\lnot}[\LogCL]
        (\pValue{v}(!B)[\LogCL])\\
      &=\pValue{v}(!A)[\LogCL].
  \end{align*}
  La première et la dernière égalité
  viennent de la \olref[val]{defn:pValue},
  la deuxième de l'hypothèse
  d'induction, et la troisième de
  \olref{prop:not}, puisque la valeur
  classique de~$!B$ appartient à
  $\{\True,\False\}$.
\item Si $!A\ident(!B\land !C)$,
  on obtient de même
  \begin{align*}
    \pValue{v}(!A)[\Log L]
      &=\tf{\land}[\Log L]
        (\pValue{v}(!B)[\Log L],
         \pValue{v}(!C)[\Log L])\\
      &=\tf{\land}[\Log L]
        (\pValue{v}(!B)[\LogCL],
         \pValue{v}(!C)[\LogCL])\\
      &=\tf{\land}[\LogCL]
        (\pValue{v}(!B)[\LogCL],
         \pValue{v}(!C)[\LogCL])\\
      &=\pValue{v}(!A)[\LogCL].
  \end{align*}
  Les étapes utilisent successivement
  la \olref[val]{defn:pValue}, l'hypothèse
  d'induction, \olref{prop:land} et
  de nouveau la définition de
  l'évaluation. Les valeurs classiques
  de~$!B$ et~$!C$ appartiennent à
  $\{\True,\False\}$.
\end{enumerate}
Les cas $!A\ident(!B\lor !C)$ et
$!A\ident(!B\lif !C)$ sont analogues,
en utilisant \olref{prop:lor} et
\olref{prop:lif}.
\end{proof}

\begin{cor}
Supposons que $\Log L$ vérifie les
hypothèses de la \olref{prop:mvl-cl},
que $\True\in V^+$ et que
$\False\notin V^+$. Pour les formules
du fragment commun défini dans cette
proposition, la conséquence de~$\Log L$
est incluse dans la conséquence
classique :
\[
  {\Entails[\Log L]}
    \subseteq {\Entails[\LogCL]}.
\]
Autrement dit, si $\Gamma$ ne contient
que des formules de ce fragment et
si~$!B$ en est une aussi, alors
$\Gamma\Entails[\Log L] !B$
implique $\Gamma\Entails[\LogCL] !B$.
En particulier, toute tautologie
de~$\Log L$ dans ce fragment est
une tautologie classique.
\end{cor}

\begin{proof}
Raisonnons par contraposition.
Supposons
$\Gamma\Entails/[\LogCL] !B$.
Il existe alors une valuation
classique
$\pAssign{v}\colon\PVar\to
\{\True,\False\}$ telle que
$\pValue{v}(!A)[\LogCL]=\True$
pour tout $!A\in\Gamma$, tandis que
$\pValue{v}(!B)[\LogCL]=\False$.
Puisque $\True,\False\in V$, cette
valuation est aussi une valuation
de~$\Log L$. Par la
\olref{prop:mvl-cl}, les valeurs
des mêmes formules dans~$\Log L$
sont respectivement $\True$ et
$\False$. Comme $\True\in V^+$
et $\False\notin V^+$, on a
$\pSat{v}{\Gamma}[\Log L]$ et
$\pSat/{v}{!B}[\Log L]$.
Donc $\Gamma\Entails/[\Log L] !B$.
\end{proof}

\begin{editorial}
La condition source écrite
$\pAssign{v}(p)\in\{\True,\False\}$
est lue ici pour chaque variable~$p$.
L'énoncé français explicite donc
les valeurs classiques de toutes
les variables utilisées. De plus,
ses hypothèses ne règlent que quatre
connecteurs : un langage peut en
contenir d'autres, notamment une
constante dont la valeur diffère
de sa valeur classique. L'énoncé
français porte donc sur le fragment
engendré par ces quatre connecteurs
et les variables, et le corollaire
est restreint au même fragment.
Dans la dernière ligne de la preuve
source, une valuation est reliée à
$\Gamma$ et~$!B$ par le symbole de
conséquence; la preuve française
emploie à cet endroit les relations
de satisfaction correspondantes.
\end{editorial}

\end{document}
